\section{密度泛函理论基本定理}

GV 第 7 章把 DFT 建立在电子液的 $E_{\mathrm{xc}}[n]$ 与耦合常数积分之上。此处给出 HK 定理、绝热连接与交换关联孔。

\subsection{Hohenberg--Kohn 定理}

\begin{theorem}[HK-1]
在非简并基态下，外势 $v_{\mathrm{ext}}(\mathbf{r})$（至多相差一个常数）由基态粒子密度 $n(\mathbf{r})$ 唯一决定。因而所有基态可观测量都是 $n$ 的泛函。
\end{theorem}

\begin{theorem}[HK-2]
存在能量泛函 $E[n]$，对给定粒子数，其最小值在真实基态密度处取得，且最小值为基态能。
\end{theorem}

\textbf{HK-1 证明提纲}（反证）：若两个相差非常数的势 $v$、$v'$ 给出同一密度 $n$，则对应基态 $|\Psi\rangle$、$|\Psi'\rangle$ 不同。由变分
\begin{equation}
  E<E'+\int(v-v')n,
  \qquad
  E'<E+\int(v'-v)n,
\end{equation}
相加得 $E+E'<E+E'$，矛盾。故势由密度唯一确定，波函数从而一切可观测量都是 $n$ 的泛函。

于是
\begin{equation}
  E[n]
  =T[n]+E_{\mathrm{H}}[n]+E_{\mathrm{xc}}[n]
  +\int v_{\mathrm{ext}}(\mathbf{r})n(\mathbf{r})\,\mathrm{d}\mathbf{r}.
\end{equation}

\subsection{绝热连接}

将相互作用标为 $\lambda v$（$0\le\lambda\le1$），同时调节外势使密度 $n(\mathbf{r})$ 保持不变。Hellmann--Feynman 给出
\begin{equation}
  E_{\mathrm{xc}}[n]
  =\frac12
  \int_0^1\mathrm{d}\lambda
  \int\mathrm{d}\mathbf{r}\,\mathrm{d}\mathbf{r}'\,
  \frac{e^2 n(\mathbf{r})n(\mathbf{r}')}{|\mathbf{r}-\mathbf{r}'|}
  \bigl[g_\lambda(\mathbf{r},\mathbf{r}')-1\bigr].
\label{eq:adiab_xc}
\end{equation}
定义耦合常数平均孔
\begin{equation}
  \bar n_{\mathrm{xc}}(\mathbf{r},\mathbf{r}')
  =n(\mathbf{r}')
  \int_0^1\mathrm{d}\lambda\,
  \bigl[g_\lambda(\mathbf{r},\mathbf{r}')-1\bigr],
\end{equation}
则
\begin{equation}
  E_{\mathrm{xc}}[n]
  =\frac12\int\mathrm{d}\mathbf{r}\,\mathrm{d}\mathbf{r}'\,
  \frac{e^2 n(\mathbf{r})\bar n_{\mathrm{xc}}(\mathbf{r},\mathbf{r}')}
  {|\mathbf{r}-\mathbf{r}'|}.
\end{equation}
$E_{\mathrm{xc}}$ 是交换关联孔的静电能。求和规则
\begin{equation}
  \int\mathrm{d}\mathbf{r}'\,\bar n_{\mathrm{xc}}(\mathbf{r},\mathbf{r}')=-1
\end{equation}
表示孔积分为 $-1$（自作用扣除）。LDA 相当于用均匀气的 $\bar n_{\mathrm{xc}}$ 就地移植。

\subsection{意义与局限}

HK 保证密度是合法变量，但不给出 $E_{\mathrm{xc}}$ 的显式。激发态需 TDDFT 等扩展。实际精度完全取决于对孔的近似。
